\newpage
\section{第15章：子类型化}

\begin{taplauthorsolutionentrybox}
\phantomsection
\label{sol:author:15.2.1}
\textbf{15.2.1 解答}

先用排列规则把 \(y\) 字段移到最前面，再用宽度规则丢掉其余字段，最后以
传递规则连接两步：
\[
\inferrule*[right=\textsc{S-Trans}]
 {\inferrule*[right=\textsc{S-RcdPerm}]{ }
   {\{x:\tapltype{Nat},y:\tapltype{Nat},z:\tapltype{Nat}\}
    \subtype
    \{y:\tapltype{Nat},x:\tapltype{Nat},z:\tapltype{Nat}\}}
  \\
  \inferrule*[right=\textsc{S-RcdWidth}]{ }
   {\{y:\tapltype{Nat},x:\tapltype{Nat},z:\tapltype{Nat}\}
    \subtype\{y:\tapltype{Nat}\}}}
 {\{x:\tapltype{Nat},y:\tapltype{Nat},z:\tapltype{Nat}\}
  \subtype\{y:\tapltype{Nat}\}}
\]

\taplsolutionbacklink{ex:15.2.1}{返回练习15.2.1}
\end{taplauthorsolutionentrybox}

\begin{taplauthorsolutionentrybox}
\phantomsection
\label{sol:author:15.2.2}
\textbf{15.2.2 解答}

不是。仍沿用正文中的缩写 \(R_x\) 与 \(R_{xy}\)。例如，可以先用包容规则
把函数 \(f\) 的类型从 \(R_x\to\tapltype{Nat}\) 改成
\(R_{xy}\to\tapltype{Nat}\)，再应用它：
\[
\inferrule*[right=\textsc{T-App}]
 {\inferrule*[right=\textsc{T-Sub}]
   {\vdash f:R_x\to\tapltype{Nat}
    \\
    \inferrule*[right=\textsc{S-Arrow}]
     {\inferrule*[right=\textsc{S-RcdWidth}]{ }{R_{xy}\subtype R_x}
      \\
      \inferrule*[right=\textsc{S-Refl}]{ }{\tapltype{Nat}\subtype\tapltype{Nat}}}
     {R_x\to\tapltype{Nat}\subtype R_{xy}\to\tapltype{Nat}}}
   {\vdash f:R_{xy}\to\tapltype{Nat}}
  \\
  \vdash xy:R_{xy}}
 {\vdash f\,xy:\tapltype{Nat}}
\]
也可以先由 \textsc{S-RcdWidth} 得到 \(R_{xy}\subtype R_x\)，再接一条
\textsc{S-Refl} 和 \textsc{S-Trans}，然后才用 \textsc{T-Sub} 把
\(xy\) 赋予 \(R_x\)。第二种做法还能任意插入更多反身和传递步骤。因此，本
演算中每个可导出的类型判断实际上都有无穷多棵不同的推导树。

\taplsolutionbacklink{ex:15.2.2}{返回练习15.2.2}
\end{taplauthorsolutionentrybox}

\begin{taplauthorsolutionentrybox}
\phantomsection
\label{sol:author:15.2.3}
\textbf{15.2.3 解答}

\begin{enumerate}
\item 共有六个：
  \[
  \{a:\tapltype{Top},b:\tapltype{Top}\},\quad
  \{b:\tapltype{Top},a:\tapltype{Top}\},\quad
  \{a:\tapltype{Top}\},\quad
  \{b:\tapltype{Top}\},\quad
  \{\},\quad
  \tapltype{Top}
  \]
\item 例如，令
  \[
  S_0=\{\},\quad S_1=\{a:\tapltype{Top}\},\quad
  S_2=\{a:\tapltype{Top},b:\tapltype{Top}\},\quad\ldots
  \]
  每增加一个字段便得到 \(S_{i+1}\subtype S_i\)。
\item 例如令 \(T_i=S_i\to\tapltype{Top}\)。箭头类型的参数位置逆变，
  所以 \(T_i\subtype T_{i+1}\)，形成无限上升链。
\end{enumerate}

\taplsolutionbacklink{ex:15.2.3}{返回练习15.2.3}
\end{taplauthorsolutionentrybox}

\begin{taplauthorsolutionentrybox}
\phantomsection
\label{sol:author:15.2.4}
\textbf{15.2.4 解答}

第一问不存在。若这样的类型存在，它只能是箭头类型或记录类型，不可能是
\(\tapltype{Top}\)。然而，记录类型不可能成为任意箭头类型的子类型，箭头
类型也不可能成为任意记录类型的子类型。

第二问同样不存在。若箭头类型 \(T_1\to T_2\) 是所有箭头类型的超类型，
根据箭头规则的逆变前提，其参数类型 \(T_1\) 必须是每个其他类型 \(S_1\)
的子类型；第一问已经说明这种类型不存在。

\taplsolutionbacklink{ex:15.2.4}{返回练习15.2.4}
\end{taplauthorsolutionentrybox}

\begin{taplauthorsolutionentrybox}
\phantomsection
\label{sol:author:15.2.5}
\textbf{15.2.5 解答}

若保留现有求值语义，加入这条规则并不安全。它会给出
\(\tapltype{Nat}\times\tapltype{Bool}\subtype\tapltype{Nat}\)，于是
本来会卡住的项
\[
  \mathsf{succ}\,(5,\tapltrue)
\]
也能通过类型检查，违反进展性定理。在\taplsecref{15.6}{sec:15.6}介绍的
强制转换语义下，这条规则是安全的；不过，即使采用该语义，它也会给子类型
检查算法带来一些困难。

\taplsolutionbacklink{ex:15.2.5}{返回练习15.2.5}
\end{taplauthorsolutionentrybox}

\begin{taplauthorsolutionentrybox}
\phantomsection
\label{sol:author:15.3.1}
\textbf{15.3.1 解答}

向子类型关系加入错误的类型对，既可能破坏保持性，也可能破坏进展性。例如，
若加入公理
\[
  \{x:\{\}\}\subtype\{x:\tapltype{Top}\to\tapltype{Top}\}
\]
便可推出 \(t=(\{x=\{\}\}.x)\{\}\) 具有类型 \(\tapltype{Top}\)。但
\(t\trans\{\}\{\}\)，而结果无法赋型，因而违反保持性。若加入公理
\[
  \{\}\subtype\tapltype{Top}\to\tapltype{Top}
\]
则项 \(\{\}\{\}\) 可以赋型，却既不是值也不能求值，因而违反进展性。

反过来，从子类型关系中删去类型对并无妨。进展性定理只在前提中提到类型关系；
限制子类型关系，继而限制类型关系，只会减少需要考虑的良类型项。保持性方面，
即使删去 \textsc{S-Trans} 也不会有问题，因为 \textsc{T-Sub} 可以恢复它
在类型推导中的作用：若已有 \(\Gamma\vdash t:S\)、\(S\subtype U\) 和
\(U\subtype T\)，可以连续使用两次 \textsc{T-Sub}，先得到
\(\Gamma\vdash t:U\)，再得到 \(\Gamma\vdash t:T\)。

\taplsolutionbacklink{ex:15.3.1}{返回练习15.3.1}
\end{taplauthorsolutionentrybox}

\begin{taplauthorsolutionentrybox}
\phantomsection
\label{sol:author:15.3.2}
\textbf{15.3.2 解答}

\begin{enumerate}
\item 对 \(S\subtype T_1\to T_2\) 的子类型推导作归纳。检查
  \cref{fig:15-1,fig:15-3}中的规则，末条规则只能是 \textsc{S-Refl}、
  \textsc{S-Trans} 或 \textsc{S-Arrow}。

  若为 \textsc{S-Refl}，则 \(S=T_1\to T_2\)，结论由反身性立即得到。
  若为 \textsc{S-Arrow}，其两个前提恰好就是所需结论。

  若为 \textsc{S-Trans}，则存在 \(U\)，使两棵子推导分别以
  \(S\subtype U\) 和 \(U\subtype T_1\to T_2\) 为结论。对第二棵子推导
  使用归纳假设，得到 \(U=U_1\to U_2\)、\(T_1\subtype U_1\) 和
  \(U_2\subtype T_2\)。现在已知 \(U\) 是箭头类型，可以对第一棵子推导
  使用归纳假设，得到 \(S=S_1\to S_2\)、\(U_1\subtype S_1\) 和
  \(S_2\subtype U_2\)。最后两次使用传递性，得到
  \(T_1\subtype S_1\) 与 \(S_2\subtype T_2\)。

\item 对 \(S\subtype\{l_i:T_i\}^{i\in1..n}\) 的子类型推导作归纳。
  末条规则只能是 \textsc{S-Refl}、\textsc{S-Trans}、
  \textsc{S-RcdWidth}、\textsc{S-RcdDepth} 或 \textsc{S-RcdPerm}。
  除传递情形外，结论都可以直接从相应规则的前提读出。

  若末条规则为 \textsc{S-Trans}，则存在记录类型
  \(U=\{u_a:U_a\}^{a\in1..o}\)，且有 \(S\subtype U\) 和
  \(U\subtype\{l_i:T_i\}^{i\in1..n}\)。对后一棵子推导使用归纳假设，
  可知每个 \(l_i\) 都是某个 \(u_a\)，且对应字段满足
  \(U_a\subtype T_i\)。再对前一棵子推导使用归纳假设，得到
  \(S=\{k_j:S_j\}^{j\in1..m}\)，每个 \(u_a\) 都是某个 \(k_j\)，且
  对应字段满足 \(S_j\subtype U_a\)。标签集合的包含关系具有传递性，
  字段类型的子类型关系也具有传递性，所以每个 \(l_i\) 都出现在 \(S\) 中，
  相应字段满足 \(S_j\subtype T_i\)。
\end{enumerate}

\taplsolutionbacklink{lem:15.3.2}{返回引理15.3.2}
\end{taplauthorsolutionentrybox}

\begin{taplauthorsolutionentrybox}
\phantomsection
\label{sol:author:15.3.6}
\textbf{15.3.6 解答}

两部分都对类型推导作归纳。这里只写第一部分；第二部分相同。

检查类型规则可知，以 \(\vdash v:T_1\to T_2\) 为结论的推导，末条规则只能
是 \textsc{T-Abs} 或 \textsc{T-Sub}。若为 \textsc{T-Abs}，结论直接
成立。若为 \textsc{T-Sub}，从前提得到 \(\vdash v:S\) 和
\(S\subtype T_1\to T_2\)。由子类型关系的反演引理，\(S\) 必为
\(S_1\to S_2\)。于是可以对前一条类型前提的推导使用归纳假设，得到
\(v\) 是 lambda 抽象。

\taplsolutionbacklink{lem:15.3.6}{返回引理15.3.6}
\end{taplauthorsolutionentrybox}
